\documentclass[a4paper]{article}
\usepackage{amsfonts}
\usepackage{amsmath}
\usepackage{amssymb}
\usepackage{graphicx}     

\begin{document}
  
\noindent \large Auteur: Odia Kibor \\
\noindent \large Studierichting: Informatica\\
\noindent \large Oefeningenreeks 1E nr. 18\\

\medskip

\normalsize


Omdat de veelterm reëel is, zijn de nulpunten: $3, 2 + i$ en $2 - i$\\

$A(z) = a((z-3)(z - (2 + i))(z - (2 - i)))$\\

Complexe paar uitwerken:\\

$(z - (2 + i))(z - (2 - i))$\\

$= ((z - 2) - i))((z - 2) + i))$\\

$= (z - 2)^2 + i(z - 2) - i(z - 2 ) + 1$\\

$= (z^2 - 4z + 4) + 1$\\

$= z^2 - 4z + 5$\\

$\Rightarrow A(z) = a((z-3)(z^2 - 4z + 5))$\\

$= a(z^3 - 4z^2 + 5z - 3z^2 + 12z -15)$\\

$= a(z^3 - 7z^2 + 17z - 15)$ met $a \in \mathbb{R}_0$







\begin{thebibliography}{999}
%\bibitem{a} Referentie
\end{thebibliography}
\end{document} 


